\scanpage{336}\mspage{252}
% l. 1–4
cela définisse un hom. de sections $F_{\alpha}\unsure{\mathbin{-}|\longrightarrow}F$. On cherche une section $G$ de $\unsure{\mathcal E|\mathscr E}$ sur $J$, et un hom de sections $F\longrightarrow G$, tel\deleted{\unsure{le}} que pour tt $\alpha$, le système \unsure{des|de} \unsure{composés|composées}
% l. 5–6
\[
 F_{\alpha}(j)\unsure{\mathbin{-}|\longrightarrow} F(j)\longrightarrow G(j)
\]
\deleted{\unsure{provienne de}} définisse un hom de sections $F_{\alpha}\longrightarrow G$.

% l. 7–11
Soit donc $\unsure{\lambda|\alpha}:j\longrightarrow j'$ une flèche de $J$, considérons les carrés \unsure{(}\uline{non commutatifs} !\unsure{,}
\diagram{\begin{tikzcd}[ampersand replacement=\&,column sep=small,row sep=small]
 \& F_{\alpha}(j)\arrow[rr,"u_{\alpha}(j)"]\arrow[dd,"F_{\alpha}(\lambda)"'] \& \& F(j)\arrow[dd,"F(\lambda)"]\\
 (\alpha)\& \& ! \&\\
 \& F_{\alpha}(j')\arrow[rr,"u_{\alpha}(j')"]\& \& F(j')
\end{tikzcd}}
(indexés par $\alpha\in A$). Pour $j'$ fixé dans $J$, le nombre de ces carrés est majoré par $\operatorname{Card}\operatorname{Ob}J_{/j'}\times\operatorname{Card}A\leq\unsure{\pi|\kappa}\mathbin{.}\pi=\pi$.

% l. 12–17
Il s'ensuit, comme $\mathcal E_{j'}$ est grande devant $\pi$, qu'il existe un objet $X(j')$ de $\mathcal E_{j'}$, et une flèche \inserted{$v(j'):$} $F(j')\unsure{\mathbin{-}|\longrightarrow} X(j')$ de telle façon que pour tout $\alpha$\inserted{$\in A$}, et tout $\lambda\in\operatorname{Ob}J_{/j'}$, le carré \deleted{\unsure{composé de la}} \inserted{\unsure{déduit|ci-dessus} de \unsure{ce|ces}}
% archive sheet number: 335
